\section{งานวิจัยและทฤษฎีที่เกี่ยวข้อง}
\label{sec:review}

\subsection{ทฤษฎีที่เกี่ยวข้อง}

\subsubsection{มาตรวัดผลลัพธ์ที่รายงานโดยผู้ป่วยและแบบประเมิน DASH และ TESS}

แบบประเมิน Disabilities of the Arm, Shoulder and Hand (DASH) พัฒนาขึ้นเพื่อวัดความบกพร่องและอาการของรยางค์บนจากมุมมองของผู้ป่วย \citep{Hudak1996} 
ส่วนแบบประเมิน Toronto Extremity Salvage Score (TESS) พัฒนาขึ้นเพื่อวัดความสามารถทางกายของผู้ป่วยเนื้องอกกระดูกและเนื้อเยื่ออ่อนภายหลังการผ่าตัดแบบสงวนรยางค์ 
โดยแยกเป็นฉบับรยางค์บนและรยางค์ล่าง \citep{Davis1996} แบบประเมินทั้งสองมีลักษณะคล้ายกัน คือข้อคำถามส่วนใหญ่ถามถึงระดับความยากลำบากในการทำกิจกรรมในชีวิตประจำวัน
และทุกข้อมีตัวเลือกคำตอบ 5 ระดับที่มีคำบรรยายกำกับ ดังรูปที่~\ref{fig:dash-tess}

\begin{figure}[htbp]
\centering
\includegraphics[width=0.6\textwidth]{figures/dash_example.png}
\caption{ตัวอย่างคำถามและตัวเลือกคำตอบของแบบประเมิน DASH ฉบับภาษาไทย}
\label{fig:dash-tess}
\end{figure}

ทั้งนี้ DASH มีข้อคำถามส่วนหนึ่งที่ถามถึงความรุนแรงของอาการและผลกระทบต่อการทำงานและกิจกรรมทางสังคม \citep{Hudak1996}
ส่วน TESS มีตัวเลือก ``ไม่เกี่ยวข้อง'' (not applicable) สำหรับกิจกรรมที่ผู้ป่วยไม่ได้ทำในชีวิตประจำวัน \citep{Davis1996}
คะแนนรายข้อ (1--5) ของทั้งสองแบบประเมินจะถูกแปลงเป็นคะแนนรวมในช่วง 0--100
โดยคะแนนของ DASH คำนวณตามสมการที่~\eqref{eq:dash_score} และต้องตอบอย่างน้อย 27 จาก 30 ข้อจึงจะคำนวณคะแนนได้ \citep{IWH2010DASHScoring}
\begin{equation}
\mathrm{DASH} = \left( \frac{1}{n} \sum_{i=1}^{n} (x_i) - 1 \right) \times 25
\label{eq:dash_score}
\end{equation}
ส่วนคะแนน TESS คำนวณตามสมการที่~\eqref{eq:tess_score} โดยไม่นำข้อที่ตอบว่าไม่เกี่ยวข้องมาคำนวณ \citep{CODE2025TESS}
\begin{equation}
\mathrm{TESS} = \frac{\sum_{i=1}^{n} (x_i) - n}{4n} \times 100
\label{eq:tess_score}
\end{equation}
โดยที่ $x_i$ คือคะแนนของข้อที่ $i$ และ $n$ คือจำนวนข้อที่ตอบ
คะแนน DASH ที่สูงหมายถึงความบกพร่องมาก ในขณะที่คะแนน TESS ที่สูงหมายถึงความสามารถทางกายที่ดี

แบบประเมินทั้งสองมีฉบับภาษาไทยที่ผ่านการตรวจสอบความตรงและความเที่ยงแล้ว \citep{Buntragulpoontawee2018, Srisawat2018} 
อย่างไรก็ตาม ผลการตรวจสอบนี้ใช้ได้กับแบบประเมินที่มีถ้อยคำและรูปแบบตามที่ใช้ในการตรวจสอบเท่านั้น
การนำแบบประเมินไปใช้ในรูปแบบอิเล็กทรอนิกส์จึงต้องคงถ้อยคำคำถามและตัวเลือกคำตอบตามต้นฉบับ
เพื่อรักษาความเท่าเทียมของการวัด (measurement equivalence) \citep{Coons2009} งานวิจัยนี้จึงออกแบบระบบให้ถามข้อคำถามด้วยถ้อยคำตามต้นฉบับทุกครั้ง

\subsubsection{การรู้จำเสียงพูดอัตโนมัติ}

การรู้จำเสียงพูดอัตโนมัติ (Automatic Speech Recognition: ASR) คือการแปลงสัญญาณเสียงพูดเป็นข้อความ แบบจำลองที่ใช้อย่างแพร่หลายในปัจจุบัน 
ได้แก่ Whisper \citep{Radford2023} ซึ่งเป็น Transformer แบบตัวเข้ารหัส-ตัวถอดรหัส (encoder-decoder) ที่รับสเปกโตรแกรม log-Mel
ของสัญญาณเสียงช่วงละ 30 วินาที และสร้างข้อความทีละโทเค็นแบบถดถอยในตัว (autoregressive) ดังรูปที่~\ref{fig:whisper_large_v3}
\begin{figure}[htbp]
\centering
\includegraphics[width=0.6\textwidth]{figures/whisper_large_v3.png}
\caption{สถาปัตยกรรมของแบบจำลอง Whisper \citep{Radford2023}}
\label{fig:whisper_large_v3}
\end{figure}

Typhoon Whisper Turbo \citep{TyphoonWhisperTurbo} ปรับจูนจาก Whisper Large-v3 Turbo ซึ่งลดจำนวนชั้นของตัวถอดรหัสของ Whisper Large-v3
จาก 32 ชั้นเหลือ 4 ชั้น ทำให้จำนวนพารามิเตอร์ลดลงจากประมาณ 1,550 ล้านเหลือ 809 ล้าน \citep{OpenAI2024WhisperTurbo}
แบบจำลองฝึกด้วยข้อมูลเสียงภาษาไทยประมาณ 11,000 ชั่วโมง ซึ่งข้อความถอดเสียงผ่านการปรับให้อยู่ในรูปคำอ่านมาตรฐาน (text normalization)
เพื่อแก้ความกำกวมของการเขียนภาษาไทย เช่น ตัวเลขที่อ่านได้หลายแบบและไม้ยมก
ซึ่งหากไม่ปรับจะทำให้ข้อมูลฝึกมีสัญญาณรบกวนและแบบจำลองทำงานได้ไม่สม่ำเสมอ \citep{Sirichotedumrong2026}

เนื่องจากภาษาไทยไม่มีการเว้นวรรคระหว่างคำ การวัดอัตราความผิดพลาดระดับคำ (Word Error Rate: WER) จึงต้องตัดคำก่อน ซึ่งทำได้ไม่ตรงไปตรงมา
\citep{Sirichotedumrong2026} จึงมีการประเมิน ASR ภาษาไทยด้วยอัตราความผิดพลาดระดับอักขระ (Character Error Rate: CER) แทน
\citep{Radford2023, Sirichotedumrong2026} ดังสมการที่~\eqref{eq:cer}
\begin{equation}
\mathrm{CER} = \frac{S + D + I}{N}
\label{eq:cer}
\end{equation}
โดยที่ $S$, $D$ และ $I$ คือจำนวนอักขระที่ถูกแทนที่ ถูกลบ และถูกแทรก ตามลำดับ และ $N$ คือจำนวนอักขระในข้อความอ้างอิง
บนชุดข้อมูลเสียงทั่วไป Gigaspeech2 และชุดข้อมูลเสียงที่มีสัญญาณรบกวน TVSpeech
Typhoon Whisper Turbo มีค่า CER ร้อยละ 4.79 และ 6.85 ตามลำดับ ใกล้เคียงกับ Typhoon Whisper Large-v3
ที่ฝึกด้วยข้อมูลชุดเดียวกัน (ร้อยละ 4.69 และ 6.32) \citep{Sirichotedumrong2026}
เนื่องด้วยขนาดที่เล็กกว่าจึงเหมาะกับการประมวลผลบนเซิร์ฟเวอร์ภายใน และแบบจำลองเหมาะกับการถอดเสียงแบบออฟไลน์
มากกว่าแบบต่อเนื่องตามเวลาจริง (streaming) \citep{TyphoonWhisperTurbo} ซึ่งสอดคล้องกับระบบที่บันทึกคำตอบของผู้ป่วยทีละข้อก่อนถอดเสียง
ทั้งนี้ ผลลัพธ์ของแบบจำลองอยู่ในรูปคำอ่าน เช่น ``สิบ'' แทน ``10'' \citep{Sirichotedumrong2026} ซึ่งส่วนตีความคำตอบต้องรองรับ


\subsubsection{แบบจำลองภาษาขนาดใหญ่ Qwen3.5}
\label{sec:qwen35}

Qwen3.5 \citep{QwenTeam2026Qwen35} เป็นแบบจำลองภาษาขนาดใหญ่แบบเปิดเผยค่าน้ำหนัก (open-weight)
ที่รับข้อมูลได้ทั้งข้อความและภาพ และรองรับ 201 ภาษา รวมถึงภาษาไทย
โดยมีแบบหนาแน่น (dense) ขนาด 0.8-27 พันล้านพารามิเตอร์ และแบบผสมผู้เชี่ยวชาญ (Mixture-of-Experts: MoE)
ขนาด 35-397 พันล้านพารามิเตอร์ จึงเลือกขนาดให้เหมาะกับเซิร์ฟเวอร์ภายในได้
และเลือกเปิดหรือปิดโหมดคิด (thinking mode) สำหรับการให้เหตุผลหลายขั้นตอนผ่านพารามิเตอร์ของแม่แบบการสนทนา \citep{QwenTeam2026Qwen35Card}

Qwen3.5 ใช้สถาปัตยกรรมของ Qwen3-Next \citep{QwenTeam2026Qwen35, QwenTeam2025Qwen3Next} ดังรูปที่~\ref{fig:qwen35-arch}
ซึ่งเป็น Transformer แบบตัวถอดรหัส (decoder-only) ที่ผสมกลไกความสนใจ 2 แบบในอัตราส่วน 3:1
แบบแรกคือ Gated DeltaNet ซึ่งเป็นความสนใจเชิงเส้น (linear attention) และแบบที่สองคือ Gated Attention ซึ่งเป็นความสนใจมาตรฐานที่เพิ่มเกตที่ผลลัพธ์ (output gate)
การผสมทั้งสองแบบให้ทั้งประสิทธิภาพและความเร็วสูงกว่าการใช้แบบใดแบบหนึ่ง
นอกจากนี้ยังใช้ Zero-Centered RMSNorm เพื่อเพิ่มเสถียรภาพของการฝึก \citep{QwenTeam2025Qwen3Next}

\begin{figure}[htbp]
\centering
\includegraphics[width=0.5\textwidth]{figures/qwen3_architecture.png}
\caption{สถาปัตยกรรมแบบผสมของ Qwen3-Next ซึ่ง Qwen3.5 ใช้เป็นฐาน \citep{QwenTeam2025Qwen3Next}}
\label{fig:qwen35-arch}
\end{figure}

% \newpage

ในงานวิจัยนี้ Qwen3.5 ถูกใช้เป็นแบบจำลองตัดสินใจ (decision model) ที่ไม่สร้างข้อความตอบกลับ
Prompt จะระบุข้อคำถาม ตัวเลือกแต่ละระดับพร้อมคำบรรยายและคำตอบของผู้ป่วย
จากนั้นแบบจำลองประมวลผลเพียงครั้งเดียวโดยที่ปิดโหมดคิด แล้วอ่านค่า logit ของโทเค็นสำหรับตัวเลือกแต่ละตัว
และคำนวณความน่าจะเป็นของแต่ละตัวเลือก ซึ่งสามารถคำนวณโดยใช้ softmax ดังสมการที่~\eqref{eq:qwen-prob}
\begin{equation}
p(o_k \mid x) = \frac{\exp(z_k)}{\sum_{j=1}^{K} \exp(z_j)}
\label{eq:qwen-prob}
\end{equation}
โดยที่ $x$ คือ prompt $z_k$ คือค่า logit ของโทเค็นที่แทนค่าของตัวเลือก $o_k$ และ $K$ คือจำนวนตัวเลือก
ตัวเลือกที่มีความน่าจะเป็นสูงสุดคือผลการจับคู่และใช้ความน่าจะเป็นนั้นเป็นค่าความมั่นใจ

% วิธีนี้ทำให้ได้ผลลัพธ์เป็นการแจกแจงความน่าจะเป็นในรูปแบบเดียวกับแบบจำลองจำแนกประเภท โดยไม่ต้องแยกวิเคราะห์ข้อความที่แบบจำลองสร้างขึ้น

% \subsubsection{แบบจำลองตัดสินใจ OpenThai-SystemOne}

% OpenThai-SystemOne \citep{iApp2026SystemOne} เป็นแบบจำลองตัดสินใจ (decision model) ภาษาไทยและภาษาอังกฤษแบบ open-weight แบบจำลองนี้ไม่มีการสร้างข้อความแต่จะรับสถานะ (state) พร้อมคำถามที่ระบุชนิด แล้วคืนค่าความน่าจะเป็นของคำตอบจากการประมวลผลเพียงครั้งเดียว

% แบบจำลองนี้ต่อยอดจาก Qwen3.5-0.8B-Base ซึ่งมีขนาดประมาณ 0.8 พันล้านพารามิเตอร์ โดยผ่านการฝึกด้วยข้อความภาษาไทยและข้อความคู่ขนานไทย-อังกฤษประมาณ 5 พันล้านโทเค็น จากนั้นแทนที่ส่วนหัวแบบจำลองภาษา (language model head) ด้วยส่วนหัวแบบช่อง (slot head) ขนาด 256 ช่อง ตัวเลือกแต่ละตัวจับคู่กับช่องหนึ่งช่อง และช่องสุดท้ายสงวนไว้สำหรับการงดตอบ (abstain)

% คำถามที่แบบจำลองรองรับมี 3 ชนิด ได้แก่ การเลือกตัวเลือก (choice) การให้คะแนนบนมาตราส่วนเชิงอันดับ (score) ซึ่งคืนค่าความน่าจะเป็นของแต่ละระดับพร้อมคะแนนเฉลี่ยถ่วงด้วยความน่าจะเป็น และคำถามแบบใช่หรือไม่ (noul) 
% ซึ่งคำถามชนิด score มีความสอดคล้องโดยตรงกับการจับคู่คำตอบของผู้ป่วยเข้ากับระดับคะแนนของข้อคำถาม และสามารถใช้ช่อง abstain ในกรณีที่ข้อมูลในคำตอบไม่เพียงพอต่อการตัดสินใจได้

\subsection{งานวิจัยที่เกี่ยวข้อง}

\subsubsection{การเตรียมชุดข้อมูลบทสนทนาสำหรับการประเมิน PROMs}
\label{sec:review-dataset}

% การพัฒนาและประเมินส่วนตีความคำตอบต้องอาศัยชุดข้อมูลคำตอบของผู้ป่วยที่มีป้ายกำกับ (label) ระดับคะแนน แต่ชุดข้อมูลลักษณะนี้มักไม่มีอยู่สำหรับแบบประเมินเฉพาะทาง อีกทั้งการเก็บข้อมูลจากผู้ป่วยจริงมีต้นทุนสูงและมีข้อจำกัดด้านความเป็นส่วนตัว งานวิจัยที่เกี่ยวข้องจึงใช้แบบจำลองภาษาขนาดใหญ่สร้างข้อมูลสังเคราะห์ (synthetic data) ทดแทน

ในพัฒนาระบบ Chat-ePRO \citep{Chen2024} ทางผู้วิจัยประสบปัญหาการขาดชุดข้อมูลในการ จึงใช้ ChatGPT สวมบทบาทแพทย์เสมือนและผู้ป่วยเสมือนสนทนากันตามข้อคำถามในแบบประเมิน ข้อคำถามแต่ละข้อนิยามเป็นเนื้อหาคำถามและชุดตัวเลือก โดยเพิ่มตัวเลือก ``ไม่แน่ใจ'' และ ``ข้าม'' เพื่อรองรับกรณีพิเศษ ผู้วิจัยสร้างบทสนทนาจากแบบประเมิน 3 ฉบับรวม 4,144 ชุด แล้วให้ผู้ตรวจสอบ 2 คนคัดกรองตามเกณฑ์ 5 ข้อ เช่น บทสนทนาต้องไม่ออกนอกเรื่อง เนื้อหาต้องไม่ขัดแย้งกันเอง และต้องไม่มีข้อมูลส่วนบุคคล คงเหลือ 3,965 ชุด และกลายเป็นชุดข้อมลูสังเคราะห์ที่ใช้สำหรับฝึกแบบจำลอง
% Tang และคณะ~\citep{Tang2023} เสนอแนวทางการใช้แบบจำลองภาษาขนาดใหญ่สร้างข้อมูลสังเคราะห์พร้อมป้ายกำกับ แล้วนำไปปรับจูนแบบจำลองขนาดเล็กที่ทำงานภายในองค์กร เพื่อหลีกเลี่ยงการส่งข้อมูลผู้ป่วยไปยังบริการภายนอกและลดภาระการเก็บและติดป้ายกำกับข้อมูล ผู้วิจัยออกแบบคำสั่งสำหรับสร้างข้อมูลแบบวนซ้ำ โดยให้ ChatGPT เสนอแม่แบบคำสั่ง 5 แบบ ทดลองสร้างข้อมูลแบบละ 10 ตัวอย่าง แล้วเลือกแม่แบบที่ให้ข้อมูลคุณภาพดีที่สุดเป็นฐานของรอบถัดไป ทำซ้ำทั้งหมด 3 รอบ ในการสร้างข้อมูลแต่ละรอบจะใส่ตัวอย่างตั้งต้น (seed) ที่สุ่มจากข้อมูลจริงลงในคำสั่งเพื่อให้ข้อมูลมีความหลากหลาย และคัดตัวอย่างที่มีคุณภาพต่ำหรือซ้ำกันออกภายหลัง เมื่อทดสอบกับงานรู้จำชื่อเฉพาะทางชีวการแพทย์ (named entity recognition) และงานสกัดความสัมพันธ์ (relation extraction) แบบจำลองที่ปรับจูนด้วยข้อมูลสังเคราะห์มีค่า F1 เฉลี่ยเพิ่มขึ้นจากร้อยละ 23.37 เป็นร้อยละ 63.99 และจากร้อยละ 75.86 เป็นร้อยละ 83.59 ตามลำดับ เมื่อเทียบกับการใช้ ChatGPT แบบ zero-shot ผลการทดลองยังพบว่าการไม่ใส่ตัวอย่างตั้งต้นทำให้ ChatGPT สร้างข้อมูลซ้ำกันและประสิทธิภาพของแบบจำลองลดลงอย่างมาก และการเพิ่มจำนวนข้อมูลสังเคราะห์ช่วยเพิ่มประสิทธิภาพได้ถึงระดับหนึ่งเท่านั้น อย่างไรก็ตาม ในงานรู้จำชื่อเฉพาะ แบบจำลองที่ปรับจูนด้วยข้อมูลสังเคราะห์ยังมีประสิทธิภาพต่ำกว่าแบบจำลองที่ปรับจูนด้วยข้อมูลจริง ซึ่งผู้วิจัยอธิบายว่าเกิดจากการกระจายตัวของข้อมูลสังเคราะห์ที่แตกต่างจากข้อมูลจริง

วิธีการที่สามารถนำมาปรับใช้ในการจัดทำชุดข้อมูลคำตอบภาษาไทยสำหรับแบบประเมิน DASH และ TESS จึงมีดังนี้
\begin{enumerate}
    \item ใช้แบบจำลองภาษาขนาดใหญ่สวมบทบาทผู้ป่วยตอบข้อคำถามทีละข้อ โดยให้ตอบด้วยการบรรยายอาการหรือวิธีทำกิจกรรมแทนการระบุตัวเลือกโดยตรง และกำหนดระดับคะแนนเป้าหมายของคำตอบไว้ล่วงหน้าเพื่อใช้เป็น label
    \item สร้างคำตอบในกรณีพิเศษ เช่น คำตอบที่กำกวมและกรณีที่ผู้ป่วยไม่ได้ทำกิจกรรมนั้น และควบคุมสัดส่วนของคำตอบแต่ละระดับให้สมดุล
    % \item ใช้ตัวอย่างคำตอบตั้งต้นจำนวนหนึ่งที่ผ่านการตรวจสอบโดยบุคลากรทางการแพทย์เป็นตัวอย่างตั้งต้นในคำสั่ง เพื่อให้ข้อมูลมีความหลากหลายและใกล้เคียงภาษาพูดของผู้ป่วย ตามข้อค้นพบของ Tang และคณะ
    \item คัดกรองข้อมูลที่สร้างขึ้นตามเกณฑ์ที่กำหนด และให้ผู้ตรวจสอบตรวจทาน label
    \item ใช้ข้อมูลสังเคราะห์สำหรับการพัฒนาและเปรียบเทียบส่วนตีความคำตอบเท่านั้น ส่วนการประเมินขั้นสุดท้ายดำเนินการกับผู้ป่วยจริง เนื่องจากข้อมูลสังเคราะห์อาจมีการกระจายตัวแตกต่างจากข้อมูลจริง
\end{enumerate}

\subsubsection{การจัดการบทสนทนา (dialogue manager) สำหรับการเก็บข้อมูล PROMs}
\label{sec:review-dialogue}

การใช้ GPT-3 ควบคุมตัวแทนสนทนาเก็บข้อมูลตามช่องข้อมูล (slot) ที่ระบุในคำสั่ง
เก็บข้อมูลได้ร้อยละ 79 โดยช่องข้อมูลท้ายคำสั่งถูกข้ามบ่อยกว่าช่องข้อมูลต้นคำสั่ง และแบบจำลองไม่ได้ถามตามมาตราส่วนที่กำหนดเสมอ ~\citep{Wei2024}
และ ChatGPT สามารถติดตามสถานะของช่องข้อมูลได้แม่นยำต่ำกว่าแบบจำลองเฉพาะงาน (joint goal accuracy 0.27 เทียบกับ 0.60) แต่เมื่อได้รับสถานะที่ถูกต้อง อัตราความสำเร็จของบทสนทนาเพิ่มจาก 0.44 เป็น 0.68 ~\citep{Hudecek2023}
ส่วน CLARITY~\citep{Shaposhnikov2025} เป็นระบบคัดกรองและส่งต่อผู้ป่วยที่ใช้เครื่องสถานะจำกัด (finite-state machine) ควบคุมบทสนทนา
และใช้แบบจำลองภาษาขนาดใหญ่สร้างข้อความภายในแต่ละ state ระบบมีการใช้งานจริงกว่า 55,000 บทสนทนา และแนะนำแพทย์เฉพาะทางได้ถูกต้องในอันดับแรกร้อยละ 77

งานวิจัยนี้จึงออกแบบให้ข้อคำถามแต่ละข้อของแบบประเมินเป็นช่องข้อมูลหนึ่งช่อง และใช้ finite-state machine ในการติดตามสถานะและควบคุมลำดับการถามแทนแบบจำลอง
รวมถึงใช้เสียงที่สร้างไว้ล่วงหน้าแทนข้อความที่แบบจำลองสร้างขึ้น เพื่อไม่ให้ข้อคำถามถูกข้ามและคงถ้อยคำตามต้นฉบับ \citep{Coons2009}

\subsubsection{แบบจำลองภาษาขนาดใหญ่สำหรับการตีความคำตอบของผู้ป่วย}
\label{sec:review-interpret}

ระบบ Chat-ePRO~\citep{Chen2024} ใช้แบบจำลอง BERT ที่ปรับจูนด้วยข้อมูลสังเคราะห์ที่ได้จากการให้ ChatGPT สวมบทบาทเป็นแพทย์และคนไข้ 
สกัดตัวเลือกจากบทสนทนา ได้ค่า Macro-F1 0.90-0.95 สูงกว่า ChatGPT แบบ few-shot ที่ได้ 0.75-0.81
HopeBot~\citep{Guo2025} ใช้ GPT-4o จัดคำตอบเข้ากับระดับของแบบประเมิน PHQ-9 และคำนวณคะแนนด้วยโปรแกรมที่แยกจากแบบจำลอง
เมื่อคำตอบกำกวมจะถามเพื่อขอความชัดเจน คะแนนจากระบบสอดคล้องกับการตอบแบบประเมินด้วยตนเองที่ค่า ICC 0.91
ส่วนการเก็บแบบสำรวจทางโทรศัพท์~\citep{Kaiyrbekov2025} ใช้ GPT-4o อนุมานคำตอบจากบทถอดความและบังคับรูปแบบของผลลัพธ์
ได้ความถูกต้องร้อยละ 98 แม้บทถอดความมีอัตราความผิดพลาดระดับคำร้อยละ 7.7
อย่างไรก็ตาม ทั้งสามงานใช้แบบจำลองที่ให้บริการผ่านระบบคลาวด์ ข้อมูลของผู้ป่วยจึงต้องถูกส่งออกนอกโรงพยาบาล
ซึ่งไม่เหมาะกับข้อมูลสุขภาพตามพระราชบัญญัติคุ้มครองข้อมูลส่วนบุคคล พ.ศ. 2562 และผู้พัฒนา Chat-ePRO ระบุว่าเป็นอุปสรรคต่อการขยายผล~\citep{Chen2024}

งานวิจัยนี้จึงเลือกใช้แบบจำลองแบบ open-weight ที่เหมาะสมกับการทำงานบนเซิร์ฟเวอร์ภายใน โดยให้แบบจำลองเป็นตัวตีความคำตอบและเลือกตัวเลือกเท่านั้น 
ส่วนการคำนวณคะแนนของแบบประเมินและการจัดการบทสนทนาทำโดยส่วนอื่นของระบบ
และกำหนดรูปแบบผลลัพธ์กลางที่มีการแจกแจงความน่าจะเป็นของทุกตัวเลือก เพื่อใช้ค่าความมั่นใจตัดสินว่าจะถามผู้ป่วยอีกครั้งหรือไม่ 
ในกรณีที่ผู้ป่วยต้องการตอบไม่ตรงตัวเลือกใดๆ หรือคำตอบไม่ชัดเจน